% Standalone problem document, generated from the adjacent README.md.
% Compile with XeLaTeX. Mathematical content is maintained in Markdown.
\documentclass[11pt,a4paper]{article}
\usepackage[fontsize=10.5pt]{scrextend}
\usepackage[margin=22mm,headheight=15pt,headsep=9mm,footskip=11mm]{geometry}
\usepackage{fontspec}
\setmainfont{texgyrepagella}[Extension=.otf,UprightFont=*-regular,BoldFont=*-bold,ItalicFont=*-italic,BoldItalicFont=*-bolditalic]
\setsansfont{texgyreheros}[Extension=.otf,UprightFont=*-regular,BoldFont=*-bold,ItalicFont=*-italic,BoldItalicFont=*-bolditalic]
\setmonofont{texgyrecursor}[Extension=.otf,UprightFont=*-regular,BoldFont=*-bold,ItalicFont=*-italic,BoldItalicFont=*-bolditalic,Scale=MatchLowercase]
\usepackage{amsmath,amssymb}
\usepackage{unicode-math}
\setmathfont{texgyrepagella-math.otf}
\usepackage{xcolor,microtype,enumitem,needspace,fancyhdr}
\usepackage{bookmark}
\definecolor{ink}{HTML}{17344B}
\definecolor{accent}{HTML}{197D80}
\definecolor{muted}{HTML}{596773}
\hypersetup{colorlinks=true,linkcolor=accent,urlcolor=accent,pdftitle={MD-01 - The
sharp Lovász-theta constant for dense random
graphs},pdfauthor={Open Problems in Numerical Linear Algebra}}
\urlstyle{same}
\setlength{\parindent}{0pt}
\setlength{\parskip}{5pt plus 1pt minus 1pt}
\linespread{1.04}
\setlength{\emergencystretch}{3em}
\widowpenalty=10000
\clubpenalty=10000
\displaywidowpenalty=10000
\allowdisplaybreaks[1]
\setlist{nosep,leftmargin=1.5em}
\providecommand{\tightlist}{\setlength{\itemsep}{0pt}\setlength{\parskip}{0pt}}
\providecommand{\pandocbounded}[1]{#1}
\pagestyle{fancy}
\fancyhf{}
\fancyhead[L]{\sffamily\footnotesize\color{muted}OPEN PROBLEMS IN NUMERICAL LINEAR ALGEBRA}
\fancyhead[R]{\sffamily\bfseries\footnotesize\color{accent}MD-01}
\fancyfoot[L]{\parbox[b]{0.9\textwidth}{\sffamily\fontsize{8}{10}\selectfont\color{muted}Editorial ratings; a bounded search cannot certify that a problem remains unsolved.\\Literature check: 2026-09-10}}
\fancyfoot[R]{\sffamily\footnotesize\color{muted}\thepage}
\renewcommand{\headrulewidth}{0.3pt}
\renewcommand{\headrule}{\hbox to\headwidth{\color{accent}\leaders\hrule height \headrulewidth\hfill}}
\makeatletter
\renewcommand{\section}{\Needspace{5\baselineskip}\@startsection{section}{1}{0pt}{12pt plus 2pt minus 2pt}{4pt}{\sffamily\bfseries\large\color{ink}}}
\renewcommand{\subsection}{\Needspace{4\baselineskip}\@startsection{subsection}{2}{0pt}{9pt plus 1pt minus 1pt}{3pt}{\sffamily\bfseries\normalsize\color{ink}}}
\makeatother
\setcounter{secnumdepth}{0}
\begin{document}
{\sffamily\small\color{muted}Matrix discrepancy and optimization\par}
\vspace{3pt}
{\raggedright\hyphenpenalty=10000\exhyphenpenalty=10000\sffamily\bfseries\fontsize{21}{24}\selectfont\color{ink}The
sharp Lovász-theta constant for dense random graphs\par}
\vspace{7pt}
{\sffamily\small\textbf{Difficulty:} challenging\quad\textbf{Importance:} broadly
interesting\par}
{\sffamily\small\bfseries\color{accent}Status: Open\par}
\vspace{4pt}
{\color{accent}\hrule height 0.8pt}
\vspace{6pt}
\textbf{Rating rationale:} A sharp expectation constant requires new
random SDP analysis beyond order bounds; it connects matrix optimization
with probabilistic combinatorics.

Let \(G_n\) be the simple random graph on \(n\) labelled vertices in
which each unordered pair is an edge independently with probability
\(1/2\). For a graph \(G\), define its Lovász number by \[
\vartheta(G)=\max\{\langle J,X\rangle:X\in\mathbb R^{n\times n},\ X=X^T\succeq0,\ \operatorname{tr}X=1,\ X_{ij}=0\text{ for }\{i,j\}\in E(G)\},
\] where \(J\) is the all-ones matrix and
\(\langle J,X\rangle=\sum_{i,j}X_{ij}\).

\textbf{Conjecture.}
\(\displaystyle\lim_{n\to\infty}\mathbb E\vartheta(G_n)/\sqrt n=1\).

This asks for the leading constant of a random semidefinite matrix
optimization problem. It is an expectation statement; convergence in
probability is not substituted for it.

\subsection{References}\label{references}

\begin{enumerate}
\def\labelenumi{\arabic{enumi}.}
\tightlist
\item
  A. S. Bandeira, D. Dmitriev, K. Lucca, P. Nizić-Nikolac, and A.
  Rödder, \emph{Randomstrasse101: Open Problems of 2025}, Entry 9,
  Conjecture 17 and the preceding SDP definition.
  \href{https://arxiv.org/abs/2603.29571}{Archived paper}.
\end{enumerate}

\subsection{Status check --- 2026-09-10}\label{status-check-2026-09-10}

The archived 2026 paper retains the conjecture. Targeted searches for
the Lovász number of dense Erdős--Rényi graphs and a sharp leading
constant through September 2026 found no resolution. The adjacent
circulant problem uses a different random matrix ensemble and is
recorded separately.

\textbf{Audit update (2026-09-10):} Rechecked the 2026 Conjecture 17 and
searched for a sharp dense-random-graph theta constant. No resolution of
the expectation limit was located. This is a bounded literature check,
not a proof that no solution exists.
\end{document}
